% GENERATED FILE. Edit canonical lessons, then run npm run generate:books.
% canonical-source-hash: fnv1a64:c832151fee35292b
% canonical-lessons: MW-C40-signs, MW-C40-notices, MW-C40-pehlo-path

\chapter{Reading — Signs, a Counter, and the Whole Bargain}
\label{ch:mw-pehlo-path}
\hlchaptermodality{40}

\begin{chapteropening}
\textbf{By the end of this chapter:} \emph{I can read Marwadi off a signboard and off a counter, and follow a whole bargain on the page without translating it word by word.}

The market exchange the track has practised out loud since chapter 13, read in order on the page for the first time, with no new word in it.
\end{chapteropening}

\section[cāy · roṭī · kapṛā · rikśā · paisā …]{(reading) — six words off a signboard}

\label{lesson:MW-C40-signs}



\begin{quote}
You have written all six of these. Not one is new.

What is new is that nobody is going to say them first.
\end{quote}

\begin{tcolorbox}[breakable,skin=enhanced,colback=orange!6,colframe=orange!45!black,boxrule=0.5pt,arc=1mm,left=6pt,right=6pt,top=4pt,bottom=4pt,fonttitle=\bfseries,title={Read this}]
\begin{quote}
\mw{चाय} \mw{रोटी} \mw{कपड़ा} \mw{रिक्शा} \mw{पैसा} \mw{भाव}
\end{quote}

Read down the list once. Do not sound out each sign — look at the whole word and let it arrive.

Again. Notice which came at once and which you had to assemble. Both are fine; the difference is what practice removes.
\end{tcolorbox}

\subsection*{What to know first}
Every one of these is a word you would actually meet painted on something — a stall, a cart, a board above a counter.

Writing a word and reading it are not the same skill. Writing gives you all the time you want and the letters in order. A signboard gives you one look, and the word has to arrive whole.

That is why six is enough for one sitting.

\subsection*{Your turn}
\begin{itemize}
\raggedright
  \item \emph{Say it:} the six words down the list, once, without stopping
\end{itemize}

\emph{Twice through:}

\begin{tcolorbox}[breakable,colback=teal!4,colframe=teal!35!black,title={Before you move on}]
How many of the six were new? (\textbf{None.})

What was new? (\textbf{Nobody said them first.} You read them.)
\end{tcolorbox}
\section[ye kitṇe kū hai? · ye ghaṇo mahaṅgo hai.]{(reading) — six lines off a counter}

\label{lesson:MW-C40-notices}



\begin{quote}
The last lesson gave you single words. These give each word something to stand on.

Still nothing new — the market lines you already say, on the page.
\end{quote}

\begin{tcolorbox}[breakable,skin=enhanced,colback=orange!6,colframe=orange!45!black,boxrule=0.5pt,arc=1mm,left=6pt,right=6pt,top=4pt,bottom=4pt,fonttitle=\bfseries,title={Read this}]
\begin{quote}
\mw{ये} \mw{कितणे} \mw{कू} \mw{है}? \mw{ये} \mw{घणो} \mw{महंगो} \mw{है।} \mw{दस} \mw{करो।} \mw{कोनी।} \mw{पैसा} \mw{ले} \mw{लो।} \mw{पाछे} \mw{मिलसू।}
\end{quote}

Read down once, without stopping.

Again — and this time hear whose line each one is. Three belong to you and three to the person behind the counter.
\end{tcolorbox}

\subsection*{What to know first}
\textbf{\mw{कोनी।}} One word, and it is a whole turn in the conversation. A notice does not have to be long to be complete.

Read in order, these six are not six notices at all — they are one bargain, from the first question to the walk-away. You have practised every one of them out loud; this is the first time you have seen them written down in the order they actually happen.

That order is the thing to carry away. A market exchange has a shape, and knowing the shape is most of reading it.

\subsection*{Your turn}
\begin{itemize}
\raggedright
  \item \emph{Say it:} the six lines, once, in order, without stopping
\end{itemize}

\emph{Twice through:}

\begin{tcolorbox}[breakable,colback=teal!4,colframe=teal!35!black,title={Before you move on}]
How many of the six lines were new? (\textbf{None.})

What did reading them in order show you? (\textbf{They are one bargain}, not six notices.)
\end{tcolorbox}
\section[ye kapṛā dikhāvo. ye kitṇe kū hai?]{(reading) — the whole bargain}

\label{lesson:MW-C40-pehlo-path}



\begin{quote}
Six words, then six lines. Now the whole thing, start to finish.

It is longer than anything you have read. Nothing in it is new.
\end{quote}

\begin{tcolorbox}[breakable,skin=enhanced,colback=orange!6,colframe=orange!45!black,boxrule=0.5pt,arc=1mm,left=6pt,right=6pt,top=4pt,bottom=4pt,fonttitle=\bfseries,title={Read this}]
\begin{quote}
\mw{ये} \mw{कपड़ा} \mw{दिखावो।} \mw{ये} \mw{कितणे} \mw{कू} \mw{है}? \mw{बीस।} \mw{ये} \mw{घणो} \mw{महंगो} \mw{है।} \mw{दस} \mw{करो।} \mw{कोनी।} \mw{पंदरा} \mw{करो।} \mw{ठीक।} \mw{पैसा} \mw{ले} \mw{लो।} \mw{पाछे} \mw{मिलसू।}
\end{quote}

Read it through once without stopping. Do not translate as you go.

Now again, and watch the number move: \mw{बीस}, \mw{दस}, \mw{पंदरा}. That is the whole negotiation, and it is carried by three words you learned as counting.
\end{tcolorbox}

\subsection*{What to know first}
\textbf{\mw{कोनी}} is the hinge. Before it, the price is falling because you pushed. After it, the price settles because they pushed back.

Notice what the passage never says. Nobody explains that \mw{पंदरा} is between \mw{दस} and \mw{बीस} — the bargain assumes you know, and you do.

This is the first time the book has asked you to hold a whole exchange at once, and the reason it can is that every line was paid for, one chapter at a time.

\subsection*{Your turn}
\begin{itemize}
\raggedright
  \item \emph{Say it:} the whole exchange aloud, once, without stopping
\end{itemize}

\emph{Twice through:}

\begin{tcolorbox}[breakable,colback=teal!4,colframe=teal!35!black,title={Before you move on}]
How many words in that exchange were new? (\textbf{None.})

Which word turned it around? (\textbf{\mw{कोनी}.})
\end{tcolorbox}
